\documentclass{IOS-Book-Article}

\usepackage{mathptmx}
\usepackage{soul}\setuldepth{article}
\usepackage{booktabs}
\usepackage{todonotes}
\usepackage{hyperref}
\usepackage{url}
% 
\def\hb{\hbox to 11.5 cm{}}

\begin{document}

\pagestyle{headings}
\def\thepage{}
\begin{frontmatter}              % The preamble begins here.


  % \pretitle{Pretitle}
  \title{LLMClean: applying OntoClean constraints through\\
    large language models to evaluate and generate taxonomies}

  \markboth{}{February 2026\hb}
  % \subtitle{Subtitle}

  \author[A]{\fnms{William} \inits{D.} \snm{Duncan}\orcid{....-....-....-....}}%
  % \thanks{Corresponding Author: Robert Hoehndorf\todo{add contact details}.}}
  \author[B]{\fnms{Robert} \snm{Hoehndorf}\orcid{....-....-....-....}}%
  % \thanks{Corresponding Author: Robert Hoehndorf\todo{add contact details}.}}
  \todo[inline]{Complete ORCID, etc.}
  
  \runningauthor{W.D. Duncan and R. Hoehndorf}
  \address[A]{\todo{Affiliation}}

  \begin{abstract}
    Taxonomies provide structure to information, yet their
    taxonomic backbones frequently suffer from misuse of the
    subsumption relation.  OntoClean, the methodology introduced
    by Guarino and Welty, addresses this problem through
    meta-properties -- rigidity, identity, unity, and dependence
    -- that constrain which ``is-a'' links are ontologically
    valid. The manual application of OntoClean is
    labor-intensive and requires ontological expertise, which
    limits its adoption.  We experimented with teaching Large
    Language Models (LLMs) the principles of OntoClean. The
    resulting method, LLMClean, automatically applies OntoClean
    constraints to evaluate and generate taxonomies.  We first
    show that LLMClean can reproduce the meta-property analysis
    from Guarino and Welty's original
    paper.
    % : Claude~4.5 Sonnet achieves 0.90~rigidity accuracy and
    % 0.62~exact match when provided with the paper text as
    % context, compared to 0.52~rigidity accuracy without it.
    We then evaluate 13 LLM on a benchmark of 10 domains (349
    terms with LLM-generated gold-standard meta-properties) in a
    zero-shot taxonomy generation task.  % Large models (GPT-4o,
    % Llama~3.3 70B, Mistral Large, Qwen~2.5 72B, DeepSeek~V3)
    % produce taxonomies with zero critical violations, while
    % small models (Llama~3.2 3B, Llama~3.1 8B) generate many.  
    We show that an agentic workflow -- pairing a Taxonomist
    agent with an OntoClean Critic agent, both backed by the
    same small model -- reduces rigidity violations from 193
    to~0 for Llama~3.2 3B and from 25 to~1 for Llama~3.1 8B.
    We further show that decomposing the critic into four
    property-specialized critics with majority voting eliminates
    all violations for Qwen~2.5 7B while retaining 86\% of the
    taxonomy links.  Our results show that LLMs can learn to
    apply OntoClean constraints and that agentic self-critique
    can recover the quality gap between small and large models.
    % \todo{revise
    % abstract after paper is complete}
  \end{abstract}

  \begin{keyword}
    OntoClean\sep ontology\sep taxonomy\sep large language models\sep
    meta-properties\sep agentic workflow
  \end{keyword}
\end{frontmatter}
\markboth{February 2026\hb}{February 2026\hb}

\section{Introduction}

\hl{The taxonomic backbone of an ontology, i.e., 
  its hierarchy of ``is-a'' (subsumption) links, is a fundamental
  structural element in ontology engineering, However, taxonomic
  relations are  also frequently
  misused~\cite{guarino2000formal}.}
\hl{Ontology engineers may confuse subsumption with part-whole, constitution,
  or role relationships, producing hierarchies in which, for example, a
  physical object is classified as an amount of matter, or where the
  ``is-a'' relation actually represents
  parthood~\cite{guarino2000formal,guarino2009ontoclean}.}
\todo{Cite: What is-a is and isn't}. 

\textcolor{blue}{Taxonomic relations (i.e., ``is-a'' (subsumption) links) form the backbone of an ontology.
  However, taxonomic relations are  also frequently misused~\cite{guarino2000formal}. Instead of subsumption, ontology engineers may misuse part-whole, constitution, or role relations. For example, a statue classified as a kind of clay confu ses the ``is-a''' relation with constitution.}

\textcolor{blue}{Such taxonomic errors} \st{These errors propagate through reasoning, integration, and reuse, and}
have motivated decades of work on principled methods for ontology
construction~\cite{guarino2009ontoclean}\todo{add more citations on ontology quality problems}.

OntoClean, hl{introduced by Guarino and}
Welty~\cite{guarino2000formal,guarino2002ontoclean,guarino2009ontoclean},
provides such a methodology.  \hl{It assigns each property (concept) in a
  taxonomy a set of \emph{meta-properties} drawn from formal ontology:
  rigidity, identity, unity, and dependence.  These meta-properties
  impose constraints on subsumption.}  \textcolor{blue}{Each taxonomic property (concept) is assigned \emph{mea-property}: rigidity, identity, unity, and dependence. The meta-properties, drawn from formal ontology, impose constraints on subsumption.} \hl{For instance, a concept assigned a rigid property
  (one that is essential to all its instances in every possible world,
  such as \emph{Person}) cannot be subsumed by a concept assigned an anti-rigid property
  (one that is contingent, such as \emph{Student}), because an entity that is necessarily a person cannot be necessarily a student.} \textcolor{blue}{For example, if \emph{Person} is assigned a rigid property, \emph{Student} is assigned an anti-rigid property, \emph{Person} cannot be subsumed by \emph{Student} because every instance of \emph{Person} is not necessarily an instance of \emph{Student}.} Violations of these constraints signal modeling errors that must be corrected before the taxonomy can be considered well-formed~\cite{guarino2000formal}.

Despite its theoretical appeal, OntoClean adoption remains limited in
practice~\cite{guarino2009ontoclean}\todo{cite evidence for limited
  adoption}.  The method requires philosophical expertise that most
domain modelers lack, and applying it manually to large ontologies is
labor-intensive: each concept must be analyzed for multiple
meta-properties, and every subsumption link must be checked against the
resulting constraints.

Large language models (LLMs) have recently shown strong performance on
a range of knowledge-intensive tasks, including biomedical ontology
matching~\cite{he2023llm4om}\todo{add more LLM-for-ontology
  citations}, taxonomy construction~\cite{babaei2023llms4ol}\todo{verify
  citation}, and conceptual classification~\todo{cite}.  Their
parametric knowledge of both philosophical concepts and domain
terminology makes them candidates for automating OntoClean analysis.
Furthermore, \emph{agentic} LLM workflows -- in which multiple
LLM-backed agents collaborate with distinct roles -- have improved
performance on complex reasoning tasks by separating generation from
critique~\cite{shinn2023reflexion}\todo{verify and add agentic LLM
  citations}.

Here, we present LLMClean, a method that uses LLMs to automatically
apply OntoClean constraints to evaluate and generate taxonomies.  We
address five research questions:

\begin{enumerate}
\item Can LLMs reproduce the meta-property analysis from the original
  OntoClean paper, and does providing the paper text as context improve
  accuracy?
\item Can LLMs generate taxonomies that respect OntoClean constraints
  in a zero-shot setting?
\item Does model scale affect the number of ontological violations?
\item Can an agentic workflow -- pairing a Taxonomist agent with an
  OntoClean Critic agent backed by the \emph{same} model -- reduce
  violations, particularly for small models?
\item Does decomposing the critic into multiple property-specialized
  critics improve taxonomy quality compared to a single monolithic
  critic?
\end{enumerate}

We evaluate 13~LLMs spanning four model families and three size
classes on a benchmark of 10~domains (349~terms).  We find that large
models (\textgreater 70B parameters or comparable proprietary models)
produce taxonomies with zero critical violations in most cases, whereas
small models (3B--9B) introduce many rigidity violations and cycles.
We then show that an agentic Taxonomist--Critic workflow eliminates
nearly all violations for small models, reducing rigidity violations
from 193 to~0 for Llama~3.2 3B.  We additionally show that
decomposing the critic into four property-specialized critics with
majority voting provides the best quality--coverage trade-off for
models at 7B parameters and above.  These results show that OntoClean
knowledge is accessible to LLMs and that structured self-critique can
recover the quality gap between small and large models.

\section{Methods}


Our method, LLMClean, builds on the OntoClean methodology in three
stages.  \textcolor{blue}{Using the OntoClean methodology, we built the LLMClean method in three stages.} First, we test whether LLMs can reproduce the meta-property analysis from Guarino and Welty's original paper -- a controlled setting with established ground truth (Section~\ref{sec:guarino-repro}).  Second, we evaluate whether LLMscan generate taxonomies that respect OntoClean constraints in a zero-shot setting across diverse domains (Section~\ref{sec:zeroshot}).  Third, we deploy an agentic workflow with a Taxonomist and an OntoClean Critic agent to improve taxonomy quality for small models (Section~\ref{sec:agentic}).  All three stages rely on shared resources and benchmark data, which we describe first.

% \hl{Our method, LLMClean, builds on the OntoClean methodology in three
%   stages.}  \textcolor{blue}{Using the OntoClean methodology, we built the LLMClean method in three stages.} First, we test whether LLMs can reproduce the meta-property analysis from Guarino and Welty's original paper -- a controlled setting with established ground truth (Section~\ref{sec:guarino-repro}).  Second, we evaluate whether LLMscan generate taxonomies that respect OntoClean constraints in a zero-shot setting across diverse domains (Section~\ref{sec:zeroshot}).  Third, we deploy an agentic workflow with a Taxonomist and an OntoClean Critic agent to improve taxonomy quality for small models (Section~\ref{sec:agentic}).  All three stages rely on shared resources and benchmark data, which we describe first.

\subsection{Background resources}

% ---------------------------------------------------------------
% Commands that produced the background resources:
% 
% 1. Download the Guarino & Welty paper:
% uv run scripts/download_paper.py \
% --url "http://cui.unige.ch/isi/cours/aftsi/articles/01-guarino00formal.pdf" \
% --output "resources/01-guarino00formal.pdf"
% 
% 2. Generate the "messy" OWL ontology:
% uv run --with rdflib scripts/generate_messy_owl.py --output ontology/guarino_messy.owl
% 
% 3. Text files for each meta-property were manually extracted from the
% converted paper text in:
% resources/converted_text_files/guarino_text_files/
% 01-guarino00formal-rigidity.txt
% 01-guarino00formal-identity.txt
% 01-guarino00formal-unity.txt
% 01-guarino00formal-dependence.txt
% 01-guarino00formal-introduction.txt
% 01-guarino00formal-converted-corrected.txt  (full paper text)
% ---------------------------------------------------------------

The OntoClean methodology requires understanding of four philosophical
notions -- identity, unity, rigidity, and dependence -- and their
formal definitions as meta-properties of
concepts~\cite{guarino2000formal}.  To make this knowledge accessible
to LLMs, we prepared text-based descriptions of each meta-property by
converting the relevant sections of Guarino and
Welty~\cite{guarino2000formal} into plain-text files.  We produced five
files, one for each meta-property (rigidity, identity, own identity,
unity, dependence), extracted from the original paper.  These files
serve as background context that can be included in LLM prompts, and
we use them in the agent-based analyzer mode (Section~\ref{sec:agentic}).

Additionally, we encoded the ground-truth classifications from the
worked examples in Guarino and Welty~\cite{guarino2000formal} as
reference data.  This ground truth includes entities such as
\emph{Person} (+R, +I, +O, +U, $-$D), \emph{Student} ($\sim$R, +I,
$-$O, +U, +D), and \emph{Red Thing} ($-$R, $-$I, $-$O, $-$U, $-$D),
and their expected OntoClean classifications (Sortal, Role, Mixin,
etc.)\todo{verify these specific ground truth values against the
  resources/ground\_truth.tsv file}.

\subsection{Benchmark dataset}
\label{sec:dataset}

We generated a benchmark dataset covering 10~diverse domains:
Treatment of bronchitis, Cycling, Naval Warfare, eCommerce,
Myrmecology, High-Frequency Trading, Urban Parkour, Quantum Computing,
Sourdough Baking, and Wastewater Treatment.  For each domain, we used
an LLM\todo{specify which model was used to generate the dataset} to
produce 35~terms (349~total across all domains, as one domain produced
34~terms) together with gold-standard meta-property assignments for
each term.

The generation prompt instructed the LLM to produce two types of terms
in each domain:

\begin{enumerate}
\item \textbf{Taxonomic core} ($\sim$40\% of terms): vertical chains of
  genuine ``is-a'' relationships with depth 3--4, forming the expected
  backbone of the taxonomy.
\item \textbf{Ontological traps} ($\sim$60\% of terms): terms related
  to the taxonomic core that are \emph{not} subtypes but are easily
  confused with them -- parts (e.g., Wing for Bird), materials (e.g.,
  Keratin), roles (e.g., Predator), and phases (e.g., Fledgling).
\end{enumerate}

Each term was annotated with four meta-properties: Rigidity ($R$:
+R, $-$R, or $\sim$R), Identity ($I$: +I or $-$I), Unity ($U$: +U,
$-$U, or $\sim$U), and Dependence ($D$: +D or $-$D), as well as a
derived ontological class (Type, Role, Phase, or Material) and a brief
explanatory note.  These meta-property assignments serve as the ground
truth for evaluating whether LLM-generated ``is-a'' links respect
OntoClean constraints.

\subsection{Models}
\label{sec:models}

We evaluated 13~LLMs from four families, spanning a range of model
sizes.  All models were accessed through the OpenRouter
API\footnote{\url{https://openrouter.ai/}}.

\paragraph{Large models.}
Gemini~2.5 Flash (Google), Claude~3.5 Sonnet (Anthropic), GPT-4o
(OpenAI), Llama~3.3 70B (Meta), Mistral Large (Mistral AI), Qwen~2.5
72B (Alibaba), and DeepSeek~V3 (DeepSeek).

\paragraph{Small models.}
Llama~3.1 8B and Llama~3.2 3B (Meta), Gemini~2.0 Flash Lite and
Gemini~2.5 Flash Lite (Google), Qwen~2.5 7B (Alibaba), and Gemma~2 9B
(Google).

For the agentic experiments, we used the small models Llama~3.2 3B,
Llama~3.1 8B, and Qwen~2.5 7B.

\section{Reproducing the Guarino and Welty taxonomy analysis}
\label{sec:guarino-repro}

% ---------------------------------------------------------------
% Commands that produced the data for this section:
% 
% 1. Generate the "messy" OWL ontology from Guarino & Welty Figure 2:
% uv run --with rdflib scripts/generate_messy_owl.py --output ontology/guarino_messy.owl
% 
% 2. Single-prompt analysis (Claude 4.5, no background):
% python3 scripts/batch_analyze_owl.py ontology/guarino_messy.owl \
% --model anthropic --output analyzed_entities_claude4-5.tsv
% 
% 3. Single-prompt analysis (Claude 4.5, with Guarino paper text):
% python3 scripts/batch_analyze_owl.py ontology/guarino_messy.owl \
% --model anthropic \
% --background-file resources/converted_text_files/guarino_text_files/01-guarino00formal-converted-corrected.txt \
% --output analyzed_entities_claude4-5_corrected_text.tsv
% 
% 4. Single-prompt analysis (Gemini 3 Flash):
% python3 scripts/batch_analyze_owl.py ontology/guarino_messy.owl \
% --model gemini --output analyzed_entities_gemini3.tsv
% 
% 5. Agent-based analysis (Claude 4.5, no background files):
% python3 scripts/batch_analyze_owl_agents.py ontology/guarino_messy.owl \
% --model anthropic --no-default-backgrounds \
% --output analyzed_entities_claude4-5_agents_no_files.tsv
% 
% 6. Agent-based analysis (Claude 4.5, with per-property background files):
% python3 scripts/batch_analyze_owl_agents.py ontology/guarino_messy.owl \
% --model anthropic \
% --output analyzed_entities_claude4-5_agents_using_files.tsv
% 
% 7. Agent-based analysis (Claude 4.5, with per-property background files + intro):
% python3 scripts/batch_analyze_owl_agents.py ontology/guarino_messy.owl \
% --model anthropic \
% --output analyzed_entities_claude4-5_agents_using_files_with_intro.tsv
% (TODO: verify exact flags for "with_intro" variant)
% 
% 8. Evaluation against ground truth:
% python3 scripts/evaluate_analysis.py analyzed_entities_claude4-5.tsv resources/ground_truth.tsv
% python3 scripts/evaluate_analysis.py analyzed_entities_claude4-5_corrected_text.tsv resources/ground_truth.tsv
% (etc. for all variants)
% ---------------------------------------------------------------

Before evaluating LLMs on generated benchmarks, we first tested
whether they can reproduce the meta-property analysis from the
original Guarino and Welty paper~\cite{guarino2000formal}.  This
experiment provides a clean evaluation setting with an established
ground truth derived directly from the paper's worked examples.

\subsection{The ``messy taxonomy'' and ground truth}

Guarino and Welty~\cite{guarino2000formal} present a ``messy
taxonomy'' (their Figure~2) as a running example of common ontological
errors.  We encoded this taxonomy as an OWL ontology
(\texttt{guarino\_messy.owl}) containing 22~classes connected by
\texttt{rdfs:subClassOf} links.  The taxonomy includes deliberate
errors: \emph{PhysicalObject} is placed as a subclass of
\emph{AmountOfMatter}, \emph{LivingBeing} is placed under
\emph{AmountOfMatter}, \emph{Person} appears under both
\emph{Vertebrate} and \emph{LegalAgent}, and \emph{Red} (an
attribute) is modeled as a class with \emph{RedApple} as its subclass.

We derived a ground truth of meta-property assignments for 21~of
these entities (excluding \emph{Entity}, the trivial root\todo{Entity
  IS in ground\_truth.tsv -- decide whether to include or exclude it
  and adjust count}) from the
paper's analysis.  Table~\ref{tab:guarino-gt} shows a representative
subset.  The full ground truth covers five meta-properties per entity:
rigidity ($R$), identity ($I$), own identity ($O$), unity ($U$), and
dependence ($D$).

\begin{table}[t]
  \caption{Representative ground-truth meta-property assignments
    derived from Guarino and Welty~\cite{guarino2000formal}.  $R$:
    rigidity, $I$: identity, $O$: own identity, $U$: unity, $D$:
    dependence.}
  \label{tab:guarino-gt}
  \centering
  \begin{tabular}{lccccc}
    \toprule
    Entity         & $R$     & $I$  & $O$  & $U$     & $D$  \\
    \midrule
    Person         & +R      & +I   & +O   & +U      & $-$D \\
    Agent          & $\sim$R & $-$I & $-$O & $-$U    & +D   \\
    Food           & $\sim$R & +I   & $-$O & $\sim$U & +D   \\
    AmountOfMatter & +R      & +I   & +O   & $\sim$U & $-$D \\
    Red            & $-$R    & $-$I & $-$O & $-$U    & $-$D \\
    RedApple       & $\sim$R & +I   & $-$O & +U      & $-$D \\
    Butterfly      & $\sim$R & +I   & $-$O & +U      & $-$D \\
    Caterpillar    & $\sim$R & +I   & $-$O & +U      & $-$D \\
    Country        & $\sim$R & +I   & +O   & +U      & +D   \\
    \bottomrule
  \end{tabular}
\end{table}

\subsection{Experimental setup}

We extracted the 22~classes from the OWL file using rdflib and
submitted each entity name (and its \texttt{rdfs:comment}, if present)
to an LLM for meta-property classification.  We tested several
configurations that vary along two dimensions: (a) the analysis
architecture (single-prompt vs.\ agent-based) and (b) the background
context provided to the LLM:

\begin{enumerate}
\item \textbf{Single-prompt, no background} (``Claude~4.5''): the LLM
  receives the entity name and a system prompt that defines the five
  meta-properties with brief definitions and examples.  The model must
  rely on its parametric knowledge to assign values.
\item \textbf{Single-prompt, with paper text} (``Claude~4.5 +
  text''): the same single-prompt setup, but the system prompt
  includes the full converted text of Guarino and
  Welty~\cite{guarino2000formal} as background context.
\item \textbf{Single-prompt, Gemini} (``Gemini~3 Flash''): the
  single-prompt setup with Gemini~3 Flash Preview, to compare across
  model families.
\item \textbf{Agent-based, no background files} (``Agents, no BG''):
  five specialized LLM agents, one per meta-property, each with a
  focused prompt but no paper excerpts.
\item \textbf{Agent-based, with per-property files} (``Agents +
  files''): five agents, each receiving the relevant section of
  Guarino and Welty (e.g., the rigidity agent receives the rigidity
  section).
\item \textbf{Agent-based, with per-property files and introduction}
  (``Agents + files + intro''): as above, but each agent also receives
  the paper's introduction for additional context\todo{verify how the
    ``with\_intro'' variant was constructed -- was the intro prepended
    to each property-specific file, or was it a separate argument?}.
\end{enumerate}

In the agent-based architecture, five separate LLM calls are made per
entity, one for each meta-property.  Each agent receives a specialized
system prompt and an optional background file containing the relevant
section of the source paper.  The own-identity agent additionally
receives the identity agent's output to enforce the constraint that
+O implies +I.  The five property values are then combined and
classified into an ontological category (Sortal, Role, Mixin, etc.)
by a deterministic rule-based classifier.

\subsection{Results: meta-property prediction accuracy}

We evaluated each configuration against the ground truth using
per-property accuracy (exact match of the predicted value against the
ground-truth value) and exact match (all five properties correct for
an entity).  Table~\ref{tab:guarino-results} summarizes the results.

\begin{table}[t]
  \caption{Per-property accuracy and exact match rate for meta-property
    prediction on the Guarino and Welty ground truth (21~entities).
    All configurations use Claude~4.5 Sonnet unless noted.  $R$:
    rigidity, $I$: identity, $O$: own identity, $U$: unity, $D$:
    dependence, EM: exact match.}
  \label{tab:guarino-results}
  \centering
  \begin{tabular}{lcccccc}
    \toprule
    Configuration & $R$  & $I$  & $O$  & $U$  & $D$  & EM   \\
    \midrule
    Claude 4.5, no background
                  & 0.52 & 0.95 & 0.71 & 0.67 & 0.67 & 0.33 \\
    Claude 4.5 + paper text
                  & 0.90 & 0.90 & 0.86 & 0.90 & 0.81 & 0.62 \\
    Gemini 3 Flash\textsuperscript{$\dagger$}
                  & 0.79 & 0.68 & 0.68 & 0.79 & 0.79 & 0.37 \\
    Agents, no BG
                  & 0.67 & 0.86 & 0.76 & 0.76 & 0.62 & 0.33 \\
    Agents + files
                  & 0.76 & 0.81 & 0.71 & 0.62 & 0.67 & 0.24 \\
    Agents + files + intro
                  & 0.71 & 0.81 & 0.76 & 0.76 & 0.57 & 0.29 \\
    \bottomrule
    \multicolumn{7}{l}{\footnotesize $\dagger$ Gemini 3 Flash evaluated on
    19/21 entities (2 JSON parse errors).}
  \end{tabular}
\end{table}

The best overall performance came from the single-prompt configuration
with the full paper text as background context (``Claude~4.5 + paper
text''), which achieved 0.90~rigidity accuracy and 0.62~exact match.
Without the paper text, the same model dropped to 0.52~rigidity
accuracy and 0.33~exact match.  This gap shows that rigidity is the
most difficult meta-property for LLMs to predict from parametric
knowledge alone -- consistent with rigidity being the most
philosophically nuanced of the five properties.

The agent-based configurations did not outperform the single-prompt
approach with paper text.  The agents-with-files configuration
(0.24~exact match) performed worse than the single-prompt baseline
(0.33~exact match).  This suggests that decomposing the analysis into
five independent calls may lose inter-property coherence: each agent
reasons about one meta-property without seeing the others, whereas a
single prompt can reason about all five jointly\todo{investigate
  whether specific entity disagreements between agents and
  single-prompt are systematic or random}.

Identity was the easiest meta-property to predict ($\geq$0.81 in five
of six configurations), while rigidity showed the largest variance
across configurations (0.52--0.90).  Dependence was consistently the
second hardest property to predict (0.57--0.81).

\subsection{Discussion of the reproduction experiment}

These results show that providing the original paper text as context
has a large positive effect, particularly for rigidity (from 0.52 to
0.90).  The LLM with paper context can distinguish anti-rigid roles
(Agent, Food, LegalAgent) from rigid types (Person, Animal,
PhysicalObject) with high accuracy, a distinction that requires
understanding the philosophical notion of essential vs.\ contingent
properties.

The finding that the agent-based approach underperformed the
single-prompt approach (with text) is notable.  We hypothesize two
reasons: (1) the single-prompt approach benefits from cross-property
reasoning (e.g., knowing that an entity is +R informs the identity
and unity decisions), and (2) each agent in the decomposed
architecture operates with less total context than the single prompt.
However, the agent-based approach may scale better: for very large
ontologies, individual property assessments can be parallelized and
are cheaper than a single complex prompt\todo{discuss whether agents
  may perform better with more sophisticated inter-agent communication,
  e.g., passing all prior property results to each subsequent agent}.

\section{Zero-shot taxonomy generation and evaluation}
\label{sec:zeroshot}

% ---------------------------------------------------------------
% How to reproduce the data for this section:
%
% Via Docker (recommended):
%   docker run --rm -e OPENROUTER_API_KEY=sk-xxx \
%     -v $(pwd)/output:/app/output llm-clean --zeroshot
%
% Via uv locally:
%
% 1. The benchmark dataset (data/benchmark_10_domains.json) is already
%    committed to the repository. To regenerate it:
%   uv run python scripts/generate_test_dataset.py data/benchmark_10_domains.json \
%     --model google/gemini-2.0-flash-001 --num-domains 10 --num-terms 35
%
% 2. Generate taxonomies for all 13 models:
%   uv run python scripts/run_benchmark.py
%   (skips models whose output file already exists)
%
%   Or run a single model:
%   uv run python scripts/generate_taxonomy.py data/benchmark_10_domains.json \
%     output/experiments/taxonomy_benchmark_gpt-4o.json --model openai/gpt-4o
%
% 3. Evaluate a single model's taxonomy:
%   uv run python scripts/evaluate_taxonomy.py data/benchmark_10_domains.json \
%     output/experiments/taxonomy_benchmark_gpt-4o.json
%
% 4. Run full benchmark and generate report:
% python3 scripts/run_benchmark.py
% (produces experiment/BENCHMARK_REPORT.md)
% ---------------------------------------------------------------

\subsection{Taxonomy generation}

We prompted each of the 13~models (Section~\ref{sec:models}) to
generate a taxonomy for each of the 10~domains in our benchmark
dataset (Section~\ref{sec:dataset}).  The prompt presented the model
with the list of terms for a given domain and instructed it to identify
strict ``is-a'' (subclass) relationships between them.  The prompt
explicitly defined the distinction between valid subsumption and common
confusions:

\begin{itemize}
\item \textbf{Valid:} ``Sparrow'' $\rightarrow$ parent ``Bird'' (a
  Sparrow \emph{is a} Bird).
\item \textbf{Invalid (part-of):} ``Wheel'' $\rightarrow$ parent
  ``Car'' (a Wheel is \emph{part of} a Car).
\item \textbf{Invalid (made-of):} ``Gold'' $\rightarrow$ parent
  ``Ring'' (a Ring is \emph{made of} Gold).
\end{itemize}

The prompt enforced a closed-world assumption: only terms from the
input list could serve as parents.  It emphasized that disconnectivity
is expected (many terms may have no parent in the list) and that
models should not force connections.  Multiple parents were permitted
when ontologically justified.  The model returned a JSON object mapping
each term to a (possibly empty) list of parent terms.

This procedure constitutes a single-pass, zero-shot generation: the
model received no examples of correct taxonomies, no OntoClean
definitions, and no feedback.  It relied entirely on its parametric
knowledge of ``is-a'' semantics.

\subsection{Evaluation metrics}

We evaluated each generated taxonomy against three OntoClean-derived
criteria using the gold-standard meta-property assignments from the
benchmark dataset:

\paragraph{Rigidity violations (critical).}
A rigid property (+R) cannot be a subclass of an anti-rigid property
($\sim$R)~\cite{guarino2000formal}.  For every parent--child link in
the generated taxonomy, we checked whether the child carries +R and
the parent carries $\sim$R.  Any such link constitutes a critical
violation.  For example, if the model places \emph{Person}~(+R) as a
subclass of \emph{Student}~($\sim$R), this is a rigidity violation.

\paragraph{Cycles (critical).}
Taxonomies must be directed acyclic graphs.  We used depth-first
search on the generated adjacency list to detect circular
dependencies (e.g., A~$\rightarrow$~B~$\rightarrow$~A).  Any cycle
constitutes a critical structural error.

\paragraph{Constitution warnings.}
An object (+I or +U) should not be classified as a subclass of a
material ($-$I and $-$U), as this typically signals a ``made-of''
relationship misrepresented as
``is-a''~\cite{guarino2000formal}.  We flagged such links as
warnings (non-critical but indicative of modeling confusion).

\subsection{Results: zero-shot benchmark}

Table~\ref{tab:benchmark} summarizes the results across all 13~models.

\begin{table}[t]
  \caption{Zero-shot taxonomy generation benchmark.  We report the total
    number of ``is-a'' links generated across all 10~domains, the number
    of critical rigidity violations, the number of cycles, and the
    number of constitution-trap warnings.  Critical violations and
    cycles must be zero for a well-formed taxonomy.}
  \label{tab:benchmark}
  \centering
  \begin{tabular}{lrrrr}
    \toprule
    Model                 & Links & Violations & Cycles & Warnings \\
    \midrule
    \multicolumn{5}{l}{\textit{Large models}}                      \\
    GPT-4o                & 193   & 0          & 0      & 1        \\
    Llama 3.3 70B         & 218   & 0          & 0      & 3        \\
    Mistral Large         & 223   & 0          & 0      & 3        \\
    Qwen 2.5 72B          & 217   & 0          & 0      & 0        \\
    DeepSeek V3           & 201   & 0          & 0      & 2        \\
    Claude 3.5 Sonnet     & 211   & 5          & 0      & 3        \\
    Gemini 2.5 Flash      & 291   & 1          & 0      & 3        \\
    \midrule
    \multicolumn{5}{l}{\textit{Small models}}                      \\
    Gemini 2.0 Flash Lite & 225   & 1          & 0      & 5        \\
    Gemini 2.5 Flash Lite & 279   & 5          & 0      & 8        \\
    Gemma 2 9B            & 261   & 7          & 2      & 7        \\
    Qwen 2.5 7B           & 375   & 14         & 31     & 4        \\
    Llama 3.1 8B          & 452   & 25         & 72     & 21       \\
    Llama 3.2 3B          & 890   & 193        & 68     & 72       \\
    \bottomrule
  \end{tabular}
\end{table}

Five large models (GPT-4o, Llama~3.3 70B, Mistral Large, Qwen~2.5
72B, and DeepSeek~V3) produced taxonomies with zero critical rigidity
violations and zero cycles across all 10~domains.  Claude~3.5 Sonnet
and Gemini~2.5 Flash, while generally strong, produced 5 and 1
violations respectively.  All large models generated between 193 and
291 links, with constitution-trap warnings remaining in the low single
digits.

Small models showed a clear quality gradient correlated with parameter
count.  Gemini~2.0 Flash Lite and Gemini~2.5 Flash Lite produced only
1 and 5 violations respectively with no cycles.  Gemma~2 9B introduced
7 violations and 2 cycles.  Qwen~2.5 7B produced 14 violations and 31
cycles.  The smallest models performed worst: Llama~3.1 8B generated
25 violations and 72 cycles, while Llama~3.2 3B produced 193 rigidity
violations, 68 cycles, and 72 constitution-trap warnings across 890
links -- an over-generation pattern in which the model connects nearly
every term to every other term without respecting ontological
constraints.

These results show that large LLMs have sufficient parametric
knowledge to respect OntoClean constraints implicitly, even without
explicit instruction.  Small models, by contrast, struggle with both
the semantics of subsumption and the structural requirement of
acyclicity.  The question that arises is whether this quality gap can be
closed without replacing the model.

\section{Agentic taxonomy generation with LLMClean}
\label{sec:agentic}

% ---------------------------------------------------------------
% How to reproduce the data for this section:
%
% === Recommended: Docker (fully self-contained, no local Python setup) ===
%
% Build the image once:
%   docker build -t llm-clean .
%
% Reproduce all agentic and multi-critic experiments
% (API key passed at runtime -- never baked into the image):
%   docker run --rm \
%     -e OPENROUTER_API_KEY=sk-xxx \
%     -v $(pwd)/output:/app/output \
%     llm-clean --agentic --multi-critic
%
% Or reproduce everything in the paper:
%   docker run --rm \
%     -e OPENROUTER_API_KEY=sk-xxx \
%     -v $(pwd)/output:/app/output \
%     llm-clean --all
%
% Using docker compose (key from host environment):
%   export OPENROUTER_API_KEY=sk-xxx
%   docker compose run --rm reproduce --agentic --multi-critic
%
% === Alternative: run locally with uv ===
%
% Step 4 -- single-critic agentic benchmark (all 3 models):
%   OPENROUTER_API_KEY=sk-xxx uv run python scripts/run_agentic_benchmark.py
%
% Or individually:
%   uv run python scripts/agentic_taxonomy.py data/benchmark_10_domains.json \
%     output/experiments/taxonomy_agentic_llama-3.2-3b-instruct.json \
%     --model meta-llama/llama-3.2-3b-instruct
%   uv run python scripts/agentic_taxonomy.py data/benchmark_10_domains.json \
%     output/experiments/taxonomy_agentic_llama-3.1-8b-instruct.json \
%     --model meta-llama/llama-3.1-8b-instruct
%   uv run python scripts/agentic_taxonomy.py data/benchmark_10_domains.json \
%     output/experiments/taxonomy_agentic_qwen-2.5-7b-instruct.json \
%     --model qwen/qwen-2.5-7b-instruct
%
% Step 5 -- multi-critic benchmark (all 3 models x 2 thresholds):
%   # majority threshold (t=2, >=2 critics must reject to veto):
%   uv run python scripts/agentic_taxonomy_multi_critic.py data/benchmark_10_domains.json \
%     output/experiments/taxonomy_agentic_multi_critic_<slug>.json \
%     --model <openrouter-model-id> --threshold 2
%
%   # strict threshold (t=1, any single critic can veto):
%   uv run python scripts/agentic_taxonomy_multi_critic.py data/benchmark_10_domains.json \
%     output/experiments/taxonomy_agentic_multi_critic_t1_<slug>.json \
%     --model <openrouter-model-id> --threshold 1
%
%   # generate comparison report across all 9 conditions:
%   uv run python scripts/run_multi_critic_benchmark.py
%   # (skips conditions whose output files already exist -- safe to resume)
%
% Evaluate any individual taxonomy file:
%   uv run python scripts/evaluate_taxonomy.py data/benchmark_10_domains.json \
%     output/experiments/taxonomy_agentic_qwen-2.5-7b-instruct.json
%
% Output files:
%   output/experiments/taxonomy_agentic_<slug>.json                  (single)
%   output/experiments/taxonomy_agentic_multi_critic_<slug>.json     (multi t=2)
%   output/experiments/taxonomy_agentic_multi_critic_t1_<slug>.json  (multi t=1)
%   output/experiments/MULTI_CRITIC_BENCHMARK_REPORT.md              (report)
% ---------------------------------------------------------------

The zero-shot results (Section~\ref{sec:zeroshot}) show that small
models produce many ontological violations when asked to generate a
complete taxonomy in a single pass.  We hypothesized that decomposing
the task into generation and critique -- with explicit OntoClean
constraints provided to the critic -- would reduce violations without
requiring a larger model.  We therefore designed an agentic workflow
with two collaborating agents, both backed by the \emph{same}
underlying LLM, and explored two critic architectures: a single
monolithic critic and a multi-critic ensemble of four specialized
property critics.

\subsection{Agent architecture}

\paragraph{Taxonomist agent.}
The Taxonomist is an LLM prompted as an ``expert Taxonomist'' whose
goal is to organize a list of terms into a strict ``is-a'' hierarchy.
It operates incrementally: in each step, it selects one unplaced term
from the input list and proposes a single parent for it (either another
term from the list or ``Thing'' if the term is a root concept).  The
Taxonomist receives the full list of terms, the current state of the
hierarchy (all previously approved links), and the set of terms
already placed.  Its prompt enforces the same constraints as the
zero-shot prompt: strict ``is-a'' only, no part-of or made-of links,
and the closed-world assumption on parent terms.

\paragraph{Single-critic architecture.}
The single critic is an LLM prompted as an ``expert in Formal
Ontology and the OntoClean methodology.''  Its system prompt includes
definitions of rigidity, identity, and unity, together with the key
OntoClean constraints:

\begin{itemize}
\item An anti-rigid property ($\sim$R) cannot subsume a rigid property
  (+R).
\item A non-sortal ($-$I) cannot subsume a sortal (+I).
\item An anti-unity property ($\sim$U) cannot subsume a unity property
  (+U).
\end{itemize}

For each link proposed by the Taxonomist, the critic infers the
meta-properties of both child and parent terms, checks the link
against these constraints, and responds with either ``APPROVE'' or
``REJECT: $\langle$explanation$\rangle$''.

\paragraph{Multi-critic architecture.}
We additionally tested a decomposed variant in which four independent
critics each evaluate a single OntoClean property:
(1)~a \emph{rigidity critic} that checks whether an anti-rigid parent
subsumes a rigid child,
(2)~an \emph{identity critic} that checks whether a non-sortal parent
subsumes a sortal child,
(3)~a \emph{unity critic} that checks whether an anti-unity parent
subsumes a unity child, and
(4)~a \emph{dependence critic} that flags cases where an independent
parent subsumes a dependent child, which can signal a role or
relational confusion.
Each critic receives a focused prompt that defines only the
meta-property it is responsible for.  We tested two rejection
thresholds: \emph{strict} ($t{=}1$), where any single critic rejection
vetoes the link, and \emph{majority} ($t{=}2$), where at least two
critics must reject before the link is vetoed.  When a link is
rejected, the aggregated feedback from all rejecting critics is passed
back to the Taxonomist for retry.

\paragraph{Interaction protocol.}
For each term in the input list (processed in alphabetical order), the
workflow proceeds as follows:

\begin{enumerate}
\item The Taxonomist proposes a link: ``\textit{child} IS-A
  \textit{parent}''.
\item The critic (single) or critics (multi) evaluate the proposed
  link against OntoClean constraints.
\item If the critic approves (or fewer than $t$ critics reject in the
  multi-critic case), the link is added to the taxonomy and the next
  term is processed.
\item If the critic rejects, the rejection reason is passed back to
  the Taxonomist, which must propose a different parent or justify its
  original choice.
\item After three failed attempts, the term is left as a root node
  (no parent assigned).
\end{enumerate}

All agents use the same LLM (same model identifier and weights) with
temperature set to 0.0 for deterministic output.  The agents share no
memory between terms: each term placement is an independent
multi-turn dialogue.  The taxonomy state (approved links) is passed as
context to the Taxonomist at each step but not to the critics, which
evaluate each link in isolation.

\subsection{Results: agentic benchmark}

We ran the agentic workflow on the same 10-domain benchmark using three
small models from the zero-shot experiment: Llama~3.2 3B, Llama~3.1
8B, and Qwen~2.5 7B.  For each model, we compared the zero-shot
baseline with the single-critic agentic workflow and the two
multi-critic thresholds ($t{=}1$ strict, $t{=}2$ majority).
Table~\ref{tab:agentic} presents the results.

\begin{table}[t]
  \caption{Zero-shot vs.\ agentic taxonomy generation for small models.
    All agentic workflows use a Taxonomist and OntoClean Critic(s)
    backed by the same model.  ``Single'' uses one monolithic critic;
    ``Multi ($t{=}2$)'' and ``Multi ($t{=}1$)'' use four
    property-specialized critics with majority and strict rejection
    thresholds, respectively.  Warn: constitution-trap warnings.}
  \label{tab:agentic}
  \centering
  \begin{tabular}{llrrrr}
    \toprule
    Model        & Mode              & Links & Viol. & Cycles & Warn. \\
    \midrule
    Llama 3.2 3B & Zero-shot         & 890   & 193   & 68     & 72    \\
                 & Single            & 160   & 0     & 6      & 1     \\
                 & Multi ($t{=}2$)   & 180   & 1     & 12     & 3     \\
                 & Multi ($t{=}1$)   & 112   & 0     & 10     & 0     \\
    \midrule
    Llama 3.1 8B & Zero-shot         & 452   & 25    & 72     & 21    \\
                 & Single            & 225   & 1     & 6      & 2     \\
                 & Multi ($t{=}2$)   & 74    & 0     & 0      & 0     \\
                 & Multi ($t{=}1$)   & 8     & 0     & 0      & 0     \\
    \midrule
    Qwen 2.5 7B  & Zero-shot         & 375   & 14    & 31     & 4     \\
                 & Single            & 255   & 3     & 15     & 3     \\
                 & Multi ($t{=}2$)   & 219   & 0     & 10     & 1     \\
                 & Multi ($t{=}1$)   & 93    & 1     & 1      & 0     \\
    \bottomrule
  \end{tabular}
\end{table}

\paragraph{Single-critic results.}
The single-critic agentic workflow produced substantial improvements
over zero-shot generation for all three models.  For Llama~3.2 3B,
rigidity violations dropped from 193 to~0 and cycles from 68 to~6.
The total number of links dropped from 890 to~160, which reflects the
critic rejecting the spurious connections that the base model produced
in zero-shot mode.  For Llama~3.1 8B, violations dropped from 25 to~1
and cycles from 72 to~6.  Qwen~2.5 7B saw violations fall from 14
to~3 and cycles from 31 to~15.

\paragraph{Multi-critic results.}
The multi-critic architecture showed different effects depending on
model capacity.  For the smallest model (Llama~3.2 3B), the
multi-critic variants did not improve over the single critic: the
majority threshold ($t{=}2$) introduced one violation and doubled the
cycle count relative to the single-critic result, while the strict
threshold ($t{=}1$) maintained zero violations but still produced more
cycles (10 vs.\ 6).  The 3B model's individual property critics appear
too unreliable for decomposed reasoning -- they approve invalid
proposals when only two of four need to agree, and they fail to catch
structural problems even when any single rejection suffices.

For the two larger models, the multi-critic majority threshold ($t{=}2$)
produced the best quality--coverage trade-off.  Llama~3.1 8B with
$t{=}2$ eliminated all violations, all cycles, and all warnings, though
at a steep coverage cost (74 links vs.\ 225 for the single critic, a
67\% reduction).  Qwen~2.5 7B with $t{=}2$ achieved the strongest
overall result: zero violations, 10 cycles (down from 15), and 1
warning, while retaining 86\% of the single-critic link count (219
vs.\ 255 links).

The strict threshold ($t{=}1$) produced near-perfect quality metrics
for both larger models but over-rejected proposals.  Llama~3.1 8B with
$t{=}1$ retained only 8 links across all 10 domains -- an essentially
empty taxonomy.  Qwen~2.5 7B with $t{=}1$ produced 93 links with just
1 violation and 1 cycle, but lost 63\% of single-critic coverage.

\paragraph{Interpretation.}
These results reveal a model-size threshold for effective critic
decomposition.  The 3B model lacks the reasoning capacity to apply
OntoClean property analysis in isolation; a single monolithic critic
that considers all constraints jointly produces better results.
Models at 7B parameters and above can serve as reliable specialized
critics, and the majority threshold ($t{=}2$) prevents individual
over-eager critics from vetoing every proposal.

The key finding remains that all agents used the \emph{same} model.
Llama~3.2 3B, a 3-billion-parameter model, possesses sufficient
knowledge of OntoClean constraints to \emph{recognize} invalid
subsumption links when explicitly prompted with constraint definitions,
even though it fails to \emph{apply} those constraints implicitly
during zero-shot generation.  The agentic architecture externalizes
the constraint-checking step that the model cannot perform internally
in a single pass.  For models at 7B and above, further decomposing
the critic into specialized property critics with majority voting
can eliminate remaining violations while preserving most of the
taxonomy coverage.

\section{Discussion}

\subsection{LLMs encode OntoClean knowledge}

Our results show that large LLMs can generate taxonomies that
respect OntoClean constraints without receiving explicit constraint
definitions in the prompt.  Five models (GPT-4o, Llama~3.3 70B,
Mistral Large, Qwen~2.5 72B, DeepSeek~V3) produced zero rigidity
violations and zero cycles across 10~domains and 349~terms.  This
suggests that the training corpora of these models contain sufficient
information about the semantics of subsumption -- and potentially
about OntoClean itself -- to avoid common ontological errors.

The constitution-trap warnings (Object ``is-a'' Material) were not
fully eliminated by any model, which indicates that the distinction
between subsumption and constitution remains difficult even for large
LLMs.  This is consistent with the observation by Guarino and
Welty~\cite{guarino2000formal} that the
constitution trap is one of the most persistent sources of
confusion in ontology modeling.

\subsection{Small models fail but can self-correct}

Small models (3B--9B parameters) produced substantially more
violations than large models, with Llama~3.2 3B generating 193
rigidity violations and 890~links -- a pattern we characterize as
``ontological hallucination,'' where the model over-generates
subsumption links without distinguishing ``is-a'' from other semantic
relations.  However, the agentic workflow showed that the same model
can recognize and reject these invalid links when given explicit
OntoClean constraints in a critic role.

The multi-critic experiment further clarified this picture.
Decomposing the critic into four property-specialized critics helped
models at 7B parameters and above: Qwen~2.5 7B with majority voting
($t{=}2$) eliminated all rigidity violations while retaining 86\% of
the taxonomy links.  For the 3B model, by contrast, the decomposed
critics were too unreliable and the single monolithic critic produced
better results.  This suggests a model-size threshold below which
decomposed ontological reasoning introduces more noise than it removes.

These findings have practical implications.  Small models are cheaper
to run, faster to deploy, and can operate in resource-constrained
environments.  The agentic workflow adds API cost (each term requires
multiple LLM calls instead of one, and the multi-critic variant
multiplies this further by a factor of four), but the total cost
remains proportional to the number of terms rather than requiring a
model upgrade.  Our results suggest that practitioners should choose
the single-critic workflow for models below 7B parameters and the
multi-critic majority variant for models at 7B and
above\todo{add cost comparison if available}.

\subsection{Limitations}

Several limitations apply to our study.

\paragraph{Gold-standard meta-properties.}
The ground-truth meta-property assignments for our benchmark dataset
were generated by an LLM rather than assigned by human ontology
experts.  While the generation prompt was designed to produce
ontologically plausible assignments, some may be
debatable\todo{validate a sample of meta-property assignments with a
  domain expert or against established ontologies}.

\paragraph{Evaluation scope.}
We evaluated only rigidity violations, cycles, and constitution
warnings.  OntoClean defines additional constraints (e.g., identity
and unity constraints on subsumption) that we did not check in our
automated evaluation.  A more thorough evaluation would verify all
OntoClean constraints\todo{consider expanding the evaluation to
  include identity and unity constraint checks}.

\paragraph{Domain coverage.}
Our benchmark covers 10~domains selected for diversity, but the
domains and terms were LLM-generated.  Performance on real-world
ontologies with hundreds or thousands of classes, ambiguous labels, and
complex axiomatization remains untested\todo{discuss plans for
  evaluation on real-world ontologies such as OBI, GO, or SNOMED CT
  fragments}.

\paragraph{Agentic cost and scalability.}
The agentic workflow processes terms sequentially and requires multiple
LLM calls per term (up to 6 in the worst case for the single critic: 3
Taxonomist proposals and 3 critic evaluations).  The multi-critic
variant increases this further: each proposal requires four critic
calls, so the worst case rises to 15 calls per term (3~proposals
$\times$ 4~critics $+$ 3~Taxonomist calls).  For large ontologies this
may become a bottleneck.  Parallelization of the four critics and
early-termination heuristics could reduce
latency\todo{quantify wall-clock time and API cost for the agentic
  experiments}.

\paragraph{Reproducibility.}
LLM outputs are non-deterministic even at temperature~0.0 for some
API providers.  We did not run repeated trials; all results represent
single runs\todo{consider reporting variance across multiple runs}.

\section{Related work}

\todo{Write related work section.  Topics to cover:
  \begin{itemize}
  \item OntoClean methodology and applications~\cite{guarino2000formal,guarino2002ontoclean,guarino2009ontoclean}
  \item Automated ontology evaluation and quality metrics
  \item LLMs for ontology engineering (ontology matching, taxonomy
    construction, concept classification): He et al.\ LLM4OM, Babaei
    Giglou et al.\ LLMs4OL, etc.
  \item Agentic LLM architectures: Reflexion, debate-based approaches
  \item Automated taxonomy construction from text (pre-LLM approaches)
  \end{itemize}
}

\section{Conclusion}

We presented LLMClean, a method that applies OntoClean constraints
through large language models to evaluate and generate taxonomies.  We
showed in three experiments that LLMs can engage with the OntoClean
methodology at different levels.  First, Claude~4.5 Sonnet reproduced
Guarino and Welty's meta-property analysis with 0.90 rigidity accuracy
when provided with the paper text -- a result that shows LLMs can
learn to apply philosophical distinctions from textual descriptions.
Second, our evaluation of 13~LLMs across 10~domains showed that large
models ($\geq$70B parameters) generate taxonomies with zero critical
ontological violations in a zero-shot setting, while small models
(3B--9B) produce many violations and cycles.  Third, an agentic
workflow in which a Taxonomist agent and an OntoClean Critic agent --
both backed by the same small model -- collaborate iteratively can
eliminate nearly all violations: rigidity violations dropped from 193
to~0 for Llama~3.2 3B.  We further showed that decomposing the critic
into four property-specialized critics with majority voting improves
the quality--coverage trade-off for models at 7B parameters and above:
Qwen~2.5 7B with the multi-critic majority variant achieved zero
violations while retaining 86\% of the single-critic link count.

These results have four implications.  First, OntoClean meta-property
analysis is accessible to LLMs, particularly when relevant background
material is provided.  Second, OntoClean constraints on subsumption
are learnable from training data by large models and can be applied
explicitly by small models in a critic role.  Third, agentic
architectures that separate generation from constraint checking offer
a practical path to high-quality ontology construction with smaller,
more accessible models.  Fourth, a model-size threshold exists for
effective critic decomposition: models below 7B parameters benefit
from a single monolithic critic, while larger models benefit from
specialized property critics with a majority rejection threshold.

Future work will extend LLMClean to evaluate existing real-world
ontologies, incorporate all OntoClean constraints (not only rigidity),
and explore whether the agentic workflow can iteratively \emph{repair}
a taxonomy rather than build one from scratch\todo{expand with
  additional future directions based on final results}.

\section*{Acknowledgements}
\todo{Add acknowledgements, funding sources.}

\section*{Data and code availability}
The source code, datasets, and experiment results are available at
\url{https://github.com/leechuck/llm-clean}\todo{verify URL is
  correct and public before submission}.

\begin{thebibliography}{99}

\bibitem{guarino2000formal}
  Guarino, N. and Welty, C. (2000).
  A Formal Ontology of Properties.
  In: \textit{Knowledge Engineering and Knowledge Management: Methods,
    Models, and Tools} (EKAW 2000), Lecture Notes in Computer Science,
  vol.~1937, pp.~97--112. Springer.

\bibitem{guarino2002ontoclean}
  Guarino, N. and Welty, C. (2002).
  Evaluating Ontological Decisions with OntoClean.
  \textit{Communications of the ACM}, 45(2):61--65.

\bibitem{guarino2009ontoclean}
  Guarino, N. and Welty, C. (2009).
  An Overview of OntoClean.
  In: Staab, S. and Studer, R. (Eds.), \textit{Handbook on Ontologies},
  2nd~ed., pp.~201--220. Springer.

\bibitem{he2023llm4om}
  \todo{Add full citation: He et al., LLM4OM or equivalent LLM
    ontology matching reference.}

\bibitem{babaei2023llms4ol}
  \todo{Add full citation: Babaei Giglou et al., LLMs4OL: Large
    Language Models for Ontology Learning.}

\bibitem{shinn2023reflexion}
  \todo{Add full citation: Shinn et al., Reflexion: Language Agents
    with Verbal Reinforcement Learning, NeurIPS 2023.}

\end{thebibliography}
\end{document}
